% Source: https://artofproblemsolving.com/community/c3h2518771_rj5303707_birthday_problems
% Stable identity: aops-topic-2518771-q07
In triangle $QRS$, the law of sines gives
\[\dfrac{RS}{\sin\angle RQS}=\dfrac{QR}{\sin\angle QSR}.\]
Since $R$ lies on line $PS$, the angle $\angle QSR$ is either equal to $\angle QSP$ (if $R$ is on ray $SP$) or supplementary to it (if $R$ is on the other side of $S$); in both cases $\sin\angle QSR=\sin\angle QSP$. In the right triangle $PQS$,
\[\sin\angle QSP=\dfrac{PQ}{PS}.\]
Substituting, with $\sin\angle RQS=\sin30^\circ=\dfrac12$:
\[2\,RS=\dfrac{QR\cdot PS}{PQ},\qquad\text{so}\qquad PS\cdot QR=2\cdot PQ\cdot RS=2\cdot7\cdot4=\boxed{56}.\]
Both positions of $R$ occur: $R$ beyond $S$ (left sketch, $\angle PQR=120^\circ$) and $R$ between $P$ and $S$ (right sketch, $\angle PQR=60^\circ$), and the argument covers both.

\begin{asy}
usepackage("amsmath");
size(4.5cm);
defaultpen(fontsize(10pt));
pair Q=(0,0), P=(7,0), Sx=(0,4.0177);
pair R=extension(P,Sx,Q,Q+dir(120));
draw(P--R,linewidth(0.8)); draw(Q--P,linewidth(0.8)); draw(Q--Sx,linewidth(0.8)); draw(Q--R,blue+linewidth(0.8));
draw((0.35,0)--(0.35,0.35)--(0,0.35));
draw(arc(Q,1.1,90,120),red); label("$30^\circ$",Q+1.1*dir(122),W,red+fontsize(7pt));
label("$Q$",Q,S); label("$P$",P,SE); label("$S$",Sx,SW); label("$R$",R,N);
label("$7$",(P+Q)/2,S); label("$4$",(R+Sx)/2,NE);
\end{asy}

\begin{asy}
usepackage("amsmath");
size(4.5cm);
defaultpen(fontsize(10pt));
pair Q=(0,0), P=(7,0), Sx=(0,7.6310);
pair R=extension(P,Sx,Q,Q+dir(60));
draw(P--Sx,linewidth(0.8)); draw(Q--P,linewidth(0.8)); draw(Q--Sx,linewidth(0.8)); draw(Q--R,blue+linewidth(0.8));
draw((0.35,0)--(0.35,0.35)--(0,0.35));
draw(arc(Q,1.4,60,90),red); label("$30^\circ$",Q+1.4*dir(75),dir(75),red+fontsize(7pt));
label("$Q$",Q,S); label("$P$",P,SE); label("$S$",Sx,N); label("$R$",R,NE);
label("$7$",(P+Q)/2,S); label("$4$",(R+Sx)/2,NE);
\end{asy}
